problem:  
Let \((M^{2n+1}, T^{1,0}M, \theta)\) be a closed strictly pseudoconvex pseudohermitian manifold whose Tanaka–Webster torsion is nonzero but divergence-free. Denote by \(\Delta_b\) its sub-Laplacian and by \(\mathrm{Ric}_b\) its pseudohermitian Ricci tensor. Suppose that  
\[
\mathrm{Ric}_b(Z, \bar Z) \ge (n+1)\kappa\,|Z|^2_\theta \quad \text{for all } Z \in T^{1,0}M
\]
for some constant \(\kappa > 0\).  
It is known that when torsion vanishes (pseudo-Einstein case), the optimal lower bound \(\lambda_1(\Delta_b) \ge n\kappa\) holds, with equality if and only if \(M\) is CR equivalent to the standard sphere (CR Obata theorem).  

**Problem:**  
Investigate whether the same rigidity conclusion remains valid when torsion is _nonvanishing but divergence-free_. More precisely, is it true that if a nontrivial pseudohermitian structure with divergence-free torsion satisfies  
\[
\lambda_1(\Delta_b) = n\kappa,
\]  
then \((M, T^{1,0}M, \theta)\) must still be CR equivalent (up to scaling) to the standard CR sphere?  
Or does there exist a new family of pseudohermitian structures with divergence-free torsion achieving the same spectral extremality but not CR equivalent to the sphere?  

Why is it a "good" problem:  
This problem builds on the known CR Lichnerowicz–Obata framework but pushes into the **torsionful regime**, a domain where classical rigidity results break down and little is understood analytically. It asks whether the presence of divergence-free Tanaka–Webster torsion can sustain the same spectral rigidity as in the torsion-free case, or whether it introduces new extremal geometries. The question is geometrically natural and simple to state—it requires only an understanding of the sub-Laplacian spectrum, Ricci curvature, and torsion—but its resolution would demand deep analytic insight into how torsion affects sub-Riemannian eigenvalue estimates. The problem therefore bridges **CR spectral geometry, pseudohermitian analysis, and contact topology**, and a solution could reveal new geometric invariants and potentially lead to a torsion-dependent extension of the CR Obata theorem.